\scanpage{94}

Notons que si $x,y$ sont des sections de $E$, on a
\[
q([x,y])=-4q(x)q(y)\unsure{\times}\underbrace{\sin^2(\widehat{x,y})}_{\substack{\Vert\,\mathrm{dfn}\\(1-\frac{\deleted{\unsure{2}}\unsure{1|2}}4\varphi(x,y)^2/q(x)q(y))}}\tag*{(21)}
\]
\margin{attention au signe et au facteur 4 !}
i.e.
\[
\deleted{4}q([x,y])=-4q(x)q(y)+\varphi(x,y)^2\tag*{(21bis)}
\]
\[
\deleted{2\varphi([x,y],[x,y])\overset{\unsure{\omega,}}{=}4(x\wedge y\wedge[x,y])}
\]
ce qui, pour $\varphi(x,y)^{\deleted{\unsure{2}}}=0$, \inserted{$q(x)=q(y)$}, donne
\[
\deleted{\unsure{\text{(22)}}\qquad\unsure{4}}q([x,y])=-4q(x)^2
\]
ou encore, en car. \inserted{\unsure{rés|résid.}} $\ne2$
\[
q\left(\deleted{\unsure{q|4}}\frac{[x,y]}{2q(x)}\right)=-1
\]
en d'autres termes on a
\[
\widetilde f:\widetilde U\longrightarrow\Sigma_1\tag*{(22)}
\]
\[
\widetilde f(\underbrace{x+iy}_{\unsure{\zeta|\xi}})=[x,y]/2q(x)
\]

Je dis que l'on a, pour $\alpha\inserted{{}=a+ib}\in\Gamma G_m^{S[i]}$
\[
\widetilde f(\alpha\zeta)=\widetilde f(\zeta).
\]

En effet
\[
\deleted{\unsure{(2)}\qquad\unsure{\widetilde f}}\ \alpha\zeta=(ax-by)+i(b\deleted{\unsure{y}}x+ay)=x'+iy'
\]
d'où
% archive sheet number: 93
